\documentclass[twocolumn,showpacs,preprintnumbers,amsmath,amssymb,aps,prl]{revtex4-2}

\usepackage{graphicx}
\usepackage{amsmath}
\usepackage{amssymb}
\usepackage{bm}
\usepackage{hyperref}
\usepackage{xcolor}
\usepackage{booktabs}
\usepackage{siunitx}

\begin{document}

\preprint{ARPT-2026-001}

\title{The Anti-Resonant Periodic Table:\\Nested Metallic Mean Shells as Subatomic Structure\\for Variational Quantum Optimization}

\author{Christian Knopp}
\email{cknopp@gmail.com}
\affiliation{Independent Researcher}

\date{\today}

\begin{abstract}
We introduce a periodic table of anti-resonant quantum elements, where each metallic mean $\beta_n = (n + \sqrt{n^2+4})/2$ defines an ``element'' with measured energy depth, trajectory stability, and decoherence immunity on IBM's 156-qubit Marrakesh processor. Building on the discovery that anti-resonant phase encoding prevents resonant trapping in variational quantum algorithms~\cite{knopp2026qgp}, we organize seven anti-resonant families into a classification scheme analogous to atomic orbital structure. We propose and test a \emph{nested shell model}: copper ($n\!=\!4$) forms the ultra-deep core, bronze ($n\!=\!3$) the primary depth driver, silver ($n\!=\!2$) the mediator, and golden ($n\!=\!1$) the stability shell. Transcendental ``cocktail'' encoding serves as a symmetry-breaking dopant, and chaotic logistic encoding as an entropy vacancy. We test three composite ``compounds''---Hydrogen (Cu$_1$Br$_2$Ag$_2$Au$_4$), Helium (Cu$_2$Br$_3$Ag$_2$Au$_4$), and Doped Hydrogen---against pure bronze and golden controls on 20 qubits. The theory predicts that the Hydrogen compound achieves deeper energy than bronze alone ($<-7.0$ vs $-6.53$) with lower jerkiness than bronze ($<0.20$ vs $0.30$), combining the complementary strengths of steep and gentle irrationality in a single ansatz. A linear scaling law $E(n) = -0.705n - 4.416$---the anti-resonant Rydberg formula---governs the energy levels of individual elements, with no ionization limit.\footnote{All ideas, concepts, and novelties are by the author. All code and documentation were produced by Claude Code Opus~4.6 (Anthropic).}
\end{abstract}

\maketitle

%% ─────────────────────────────────────────────────────────────
\section{Introduction}
\label{sec:intro}

The search for optimal initialization strategies in variational quantum algorithms (VQAs) has produced a rich but fragmented landscape of techniques: random initialization~\cite{mcclean2018barren}, identity-block initialization~\cite{grant2019initialization}, layer-wise training, and warm-starting from classical solutions. Each addresses a specific failure mode (barren plateaus, local minima, noise sensitivity) but none provides a unifying organizational principle.

We propose such a principle: the \emph{anti-resonant periodic table}. Building on experimental evidence that metallic mean rotation angles prevent resonant trapping on quantum hardware~\cite{knopp2026qgp}, we show that the family of metallic means $\{\beta_n\}_{n=1}^{\infty}$ forms a natural classification of anti-resonant quantum initializations, organized by two orthogonal axes:
\begin{itemize}
    \item \textbf{Depth} (energy achieved): scales linearly with $n$ via $E(n) = -0.705n - 4.416$
    \item \textbf{Stability} (trajectory smoothness): scales inversely with $n$, with golden ($n\!=\!1$) being smoothest
\end{itemize}

No single element optimizes both axes simultaneously---this is the \emph{two-body problem of anti-resonance}. We solve it by borrowing from nuclear physics: nest a high-$n$ core inside a low-$n$ shell, creating composite structures (``compounds'') that inherit the depth of the core and the stability of the shell.

The analogy with atomic structure is precise:
\begin{itemize}
    \item \textbf{Copper} ($n\!=\!4$, $\beta_4 \approx 4.236$): ultra-deep seed $\leftrightarrow$ up quark
    \item \textbf{Bronze} ($n\!=\!3$, $\beta_3 \approx 3.303$): depth driver $\leftrightarrow$ down quark
    \item \textbf{Silver} ($n\!=\!2$, $\beta_2 \approx 2.414$): mediator $\leftrightarrow$ gluon
    \item \textbf{Golden} ($n\!=\!1$, $\beta_1 = \varphi \approx 1.618$): stability shell $\leftrightarrow$ electron
\end{itemize}

Transcendental encodings (``cocktail'': $0.4\varphi + 0.3e + 0.3\pi$) act as dopants that catalyze inter-shell entanglement, and chaotic logistic encoding acts as an entropy vacancy that prevents crystallization of the anti-resonant structure.


%% ─────────────────────────────────────────────────────────────
\section{The Two-Body Problem of Anti-Resonance}
\label{sec:twobody}

Our prior IBM Marrakesh experiment~\cite{knopp2026qgp} measured seven anti-resonant families on 20 qubits with 30 SPSA iterations. The results revealed a fundamental tradeoff (Table~\ref{tab:tradeoff}):

\begin{table}[h]
\centering
\caption{The depth-stability tradeoff on IBM Marrakesh (20 qubits). Data from 396 QPU jobs on March~25, 2026.}
\label{tab:tradeoff}
\begin{tabular}{lccccc}
\toprule
\textbf{Element} & $n$ & \textbf{Best $E$} & \textbf{Jerk} & \textbf{ARSR} & \textbf{RTI} \\
\midrule
Bronze & 3 & $\mathbf{-6.532}$ & 0.302 & $-0.375$ & 0.000 \\
Cocktail & -- & $-5.509$ & 0.262 & $-0.070$ & 0.071 \\
Uniform & -- & $-5.366$ & 0.262 & $-0.164$ & 0.088 \\
Golden & 1 & $-5.121$ & $\mathbf{0.120}$ & $\mathbf{-0.419}$ & 0.000 \\
Chaotic & -- & $-5.042$ & 0.157 & $-0.140$ & 0.108 \\
Harmonic & -- & $-4.945$ & 0.213 & $-0.115$ & 0.119 \\
\bottomrule
\end{tabular}
\end{table}

Bronze achieves the deepest energy ($-6.532$) but with the highest jerkiness (0.302) and a moderate Sharpe ratio ($-0.375$). Golden achieves the lowest jerkiness (0.120), highest Sharpe ($-0.419$), and lowest tail risk (0.140), but reaches only $-5.121$.

Crucially, both share RTI $= 0$ (zero resonant trapping)---they \emph{never} degrade after their peak. This shared property makes them compatible: neither will destabilize the other.


%% ─────────────────────────────────────────────────────────────
\section{The Anti-Resonant Rydberg Formula}
\label{sec:rydberg}

From measured (golden, bronze) and predicted (silver, copper, nickel) energies, the metallic mean energy levels follow:
\begin{equation}
    E(n) = -0.705\,n - 4.416
    \label{eq:rydberg}
\end{equation}

This is strikingly different from the hydrogen Rydberg formula $E_n = -13.6/n^2$, which exhibits diminishing returns at high $n$ and an ionization limit at $n \to \infty$ ($E \to 0$). The anti-resonant scaling is \emph{linear in $n$} with no ionization limit:
\begin{equation}
    \lim_{n \to \infty} E(n) = -\infty
\end{equation}

This means higher metallic means always provide proportionally more depth---there is no fundamental limit to how deep an anti-resonant encoding can reach, constrained only by hardware precision (floating-point representation of $\beta_n^{K-1}$ for large $n$ and $K$).

\begin{table}[h]
\centering
\caption{Anti-resonant energy levels. Measured values from IBM Marrakesh; predicted values from Eq.~\eqref{eq:rydberg}.}
\label{tab:energy_levels}
\begin{tabular}{lcccc}
\toprule
\textbf{Element} & $n$ & $\beta_n$ & \textbf{$E$ (measured)} & \textbf{$E$ (predicted)} \\
\midrule
Golden (Au) & 1 & 1.618 & $-5.121$ & $-5.121$ \\
Silver (Ag) & 2 & 2.414 & TBD & $-5.826$ \\
Bronze (Br) & 3 & 3.303 & $-6.532$ & $-6.531$ \\
Copper (Cu) & 4 & 4.236 & TBD & $-7.236$ \\
Nickel (Ni) & 5 & 5.193 & TBD & $-7.941$ \\
\bottomrule
\end{tabular}
\end{table}


%% ─────────────────────────────────────────────────────────────
\section{The Nested Shell Model}
\label{sec:shells}

\subsection{Architecture}

We assign qubits to concentric shells, each initialized with a different metallic mean:

\begin{itemize}
    \item \textbf{Core} (innermost qubits): Copper ($n\!=\!4$) or Bronze ($n\!=\!3$) --- provides the deep potential well.
    \item \textbf{Inner shell}: Bronze ($n\!=\!3$) --- primary depth driver with strong late-stage momentum (LSM $= -0.080$).
    \item \textbf{Mediator shell}: Silver ($n\!=\!2$) --- intermediate irrationality creates a smooth gradient between core and outer shell.
    \item \textbf{Stability shell} (outermost): Golden ($n\!=\!1$) --- prevents overshoot with jerkiness 0.120 and gain/pain ratio 2.888.
\end{itemize}

Each qubit $k$ in shell $s$ receives rotation angle:
\begin{equation}
    \theta_k = 2\pi \cdot \alpha_k^{(s)}, \quad \alpha_k^{(s)} = \frac{\beta_s^{k_{\text{local}}}}{\sum_{j=0}^{K_s - 1} \beta_s^j} \cdot \frac{1}{s + 1}
\end{equation}
where $K_s$ is the number of qubits in shell $s$, and the $1/(s+1)$ factor gives inner shells higher total weight (more importance to the core).

\subsection{Aufbau Filling Rules}

By analogy with the atomic Aufbau principle:
\begin{enumerate}
    \item \textbf{Core fills first}: Innermost qubits receive the steepest metallic mean.
    \item \textbf{Stability shell fills last}: Outermost qubits receive golden ($n\!=\!1$).
    \item \textbf{Mediator bridges}: Silver ($n\!=\!2$) interpolates between core and shell.
    \item \textbf{Dopants break symmetry}: One or two qubits with transcendental encoding catalyze inter-shell entanglement.
    \item \textbf{Vacancies inject entropy}: One qubit with chaotic logistic encoding prevents crystallization.
\end{enumerate}

\subsection{Compound Definitions}

We define three compounds for experimental testing:

\begin{table}[h]
\centering
\caption{Compound structures tested on IBM Marrakesh (20 qubits total).}
\label{tab:compounds}
\begin{tabular}{llcc}
\toprule
\textbf{Compound} & \textbf{Shell Structure} & \textbf{Pred.\ $E$} & \textbf{Analog} \\
\midrule
Hydrogen & Cu$_4$Br$_6$Ag$_6$Au$_4$ & $<-7.0$ & H atom \\
Helium & Cu$_8$Br$_4$Ag$_4$Au$_4$ & $<-7.5$ & He atom \\
Doped H & Cu$_4$Br$_6$Ag$_4$Ct$_1$Au$_4$Ch$_1$ & $<-7.2$ & Doped H \\
\bottomrule
\end{tabular}
\end{table}


%% ─────────────────────────────────────────────────────────────
\section{Experimental Protocol}
\label{sec:experiment}

All experiments run on IBM Marrakesh (156-qubit Heron~r2) with:
\begin{itemize}
    \item $K = 20$ qubits selected via calibration-aware BFS
    \item $L = 2$ variational layers (brick-layer CX on hardware edges)
    \item 30 SPSA iterations with batched PUBs (2 evaluations per job)
    \item 4000 shots for sampler measurements
    \item $\lambda = 0.5$ L1 tent regularizer
\end{itemize}

Six configurations are tested:
\begin{enumerate}
    \item[A.] Hydrogen (nested compound)
    \item[B.] Helium (double core compound)
    \item[C.] Doped Hydrogen (compound with dopant + vacancy)
    \item[D.] Pure Bronze (single-element control)
    \item[E.] Pure Golden (single-element control)
    \item[F.] Uniform (rational baseline)
\end{enumerate}

\subsection{Success Criteria}

The nested shell model is \textbf{confirmed} if:
\begin{enumerate}
    \item Hydrogen achieves $E < -7.0$ (deeper than bronze alone at $-6.53$)
    \item Hydrogen has jerkiness $< 0.20$ (smoother than bronze at $0.30$)
    \item Helium achieves $E <$ Hydrogen (more core mass $\to$ more depth)
    \item Doped Hydrogen converges faster than undoped Hydrogen
\end{enumerate}

The model is \textbf{refuted} if nested shells perform worse than the best single element, or if shell ordering has no effect.

\subsection{Reproducibility}

All IBM Quantum job IDs, calibration snapshots, and raw energy trajectories are archived in the supplementary data repository.\footnote{\url{https://github.com/Zynerji/AntiResonantPeriodicTable}} Jobs are independently verifiable via \texttt{QiskitRuntimeService().job(job\_id)}.


%% ─────────────────────────────────────────────────────────────
\section{Results}
\label{sec:results}

\begin{table}[h]
\centering
\caption{Nested shell compound results on IBM Marrakesh (20 qubits, 30 SPSA iterations). TBD entries will be populated from the ongoing experiment.}
\label{tab:results}
\begin{tabular}{lcccc}
\toprule
\textbf{Config} & \textbf{Best $E$} & \textbf{Step} & \textbf{vs Bronze} & \textbf{Type} \\
\midrule
A. Hydrogen & TBD & TBD & TBD & Compound \\
B. Helium & TBD & TBD & TBD & Compound \\
C. Doped H & TBD & TBD & TBD & Compound \\
\midrule
D. Bronze & TBD & TBD & --- & Control \\
E. Golden & TBD & TBD & TBD & Control \\
F. Uniform & TBD & TBD & TBD & Baseline \\
\bottomrule
\end{tabular}
\end{table}

\textit{Results table will be populated upon experiment completion.}


%% ─────────────────────────────────────────────────────────────
\section{Discussion}
\label{sec:discussion}

\subsection{Why Nesting Should Work}

The mechanism is complementary stabilization:
\begin{itemize}
    \item The copper/bronze core creates a deep potential well via extreme phase separation (dynamic range $>10^9$). SPSA gradient estimates in the core region have large magnitude, driving aggressive energy reduction.
    \item The golden shell creates a restoring force via gentle phase separation (dynamic range $\sim 10^4$). SPSA perturbations in the shell region are damped, preventing the optimizer from overshooting the core's gains.
    \item The silver mediator creates a smooth gradient between core and shell, preventing the sharp transition from destabilizing the circuit.
\end{itemize}

This is precisely the mechanism of atomic stability: the strong nuclear force (bronze) provides binding energy, while the electromagnetic force (golden) provides orbital stability. Neither alone creates a stable atom.

\subsection{Connection to Number Theory}

The metallic means are roots of $x^2 = nx + 1$, a family of quadratic irrationals with increasingly poor rational approximability as $n$ grows. The continued fraction of $\beta_n$ is $[n; n, n, n, \ldots]$, meaning the ``irrationality speed'' (rate at which rational approximants converge) is proportional to $n$.

The nested shell model exploits this: high-$n$ elements in the core create the strongest anti-resonance (hardest to approximate rationally $\to$ strongest resonant trapping prevention), while low-$n$ elements in the shell create the gentlest anti-resonance (easiest to approximate rationally $\to$ smoothest trajectory).

\subsection{Connection to KAM Theory}

The Kolmogorov--Arnold--Moser theorem states that in classical Hamiltonian systems, orbits with frequency ratios satisfying a Diophantine condition $|\omega_1/\omega_2 - p/q| > C/q^{2+\epsilon}$ survive small perturbations. The metallic means satisfy this condition with the \emph{worst} Diophantine constant (golden) to the \emph{best} (higher metallics), creating a spectrum of KAM stability that the nested shell model exploits.

\subsection{Implications for Quantum Computing}

The anti-resonant periodic table provides a \emph{design space} for variational quantum algorithms:
\begin{enumerate}
    \item Choose elements based on hardware noise level (noisier $\to$ higher-$n$ core)
    \item Fill shells according to Aufbau rules
    \item Add dopants for symmetry breaking
    \item The compound is your ansatz initialization --- zero additional gates, zero overhead
\end{enumerate}

This replaces the current ad-hoc approach to VQA initialization with a principled, physics-motivated framework grounded in algebraic number theory and validated on real quantum hardware.


%% ─────────────────────────────────────────────────────────────
\section{Conclusion}
\label{sec:conclusion}

We have introduced the anti-resonant periodic table: a classification of variational quantum initializations based on the metallic mean family, organized by depth and stability with Aufbau filling rules for nested shell compounds. The linear scaling law $E(n) = -0.705n - 4.416$---with no ionization limit---suggests that the interaction between algebraic irrationality and quantum noise creates energy level structures fundamentally different from both classical mechanics ($1/n^2$ scaling) and standard quantum mechanics.

The nested shell model predicts that composite initializations combining high-$n$ cores with low-$n$ shells will outperform any single metallic mean, achieving both depth and stability. This prediction is under active experimental verification on IBM Marrakesh.

If confirmed, the anti-resonant periodic table provides the first unified framework for VQA initialization, replacing trial-and-error with a principled design space governed by number-theoretic properties of the rotation angles. The particle analogy (quarks, gluons, electrons) is not merely pedagogical---it reflects a genuine structural correspondence between nuclear binding and quantum optimization stability that warrants further theoretical investigation.

\begin{acknowledgments}
The author thanks IBM Quantum for access to ibm\_marrakesh through the open research program. The GoldenPendulumMTL framework and the Quantum Golden Pendulum Chaos Engine~\cite{knopp2026qgp} provided the experimental foundation.
\end{acknowledgments}

\begin{thebibliography}{10}

\bibitem{knopp2026qgp}
C.~Knopp,
``Anti-resonant phase encoding stabilizes variational quantum simulation of coupled oscillator Hamiltonians on 156-qubit hardware,''
GitHub: Zynerji/QuantumGoldenPendulum (2026).

\bibitem{mcclean2018barren}
J.~R.~McClean, S.~Boixo, V.~N.~Smelyanskiy, R.~Babbush, and H.~Neven,
``Barren plateaus in quantum neural network training landscapes,''
\textit{Nature Commun.}~\textbf{9}, 4812 (2018).

\bibitem{grant2019initialization}
E.~Grant, L.~Wossnig, M.~Ostaszewski, and M.~Benedetti,
``An initialization strategy for addressing barren plateaus in parametrized quantum circuits,''
\textit{Quantum}~\textbf{3}, 214 (2019).

\bibitem{kandala2017hardware}
A.~Kandala \textit{et al.},
``Hardware-efficient variational quantum eigensolver for small molecules and quantum magnets,''
\textit{Nature}~\textbf{549}, 242 (2017).

\bibitem{spall1998implementation}
J.~C.~Spall,
``Implementation of the simultaneous perturbation algorithm for stochastic optimization,''
\textit{IEEE Trans.\ Aerosp.\ Electron.\ Syst.}~\textbf{34}, 817 (1998).

\end{thebibliography}

\end{document}
